% !TEX program = xelatex
\documentclass[11pt,fontset=fandol]{ctexart}

\usepackage[a4paper,margin=2.0cm]{geometry}
\usepackage{booktabs}
\usepackage{tabularx}
\usepackage{longtable}
\usepackage{array}
\usepackage{xcolor}
\usepackage{hyperref}
\usepackage{enumitem}

\definecolor{linkblue}{HTML}{2F6FAB}
\definecolor{muted}{HTML}{5D6D7E}
\definecolor{softred}{HTML}{B03A2E}
\definecolor{softgreen}{HTML}{1E8449}
\definecolor{softorange}{HTML}{A65F00}

\hypersetup{
  colorlinks=true,
  linkcolor=linkblue,
  urlcolor=linkblue
}

\setlength{\parindent}{0pt}
\setlength{\parskip}{0.55em}
\emergencystretch=3em
\renewcommand{\arraystretch}{1.22}
\newcolumntype{Y}{>{\raggedright\arraybackslash}X}
\newcolumntype{L}[1]{>{\raggedright\arraybackslash}p{#1}}
\newcommand{\code}[1]{\texttt{\detokenize{#1}}}
\newcommand{\good}[1]{\textcolor{softgreen}{#1}}
\newcommand{\bad}[1]{\textcolor{softred}{#1}}
\newcommand{\warn}[1]{\textcolor{softorange}{#1}}
\urlstyle{same}

\title{\bfseries OpenCollab SWE-bench V3 评测 Harness 说明与审核报告}
\author{Kai / Codex}
\date{2026-07-06}

\begin{document}
\sloppy
\maketitle

\section{目标}

OpenCollab 的 SWE-bench 做题框架把四件事合并管理：生成 patch、保存可恢复现场、发现可评测 patch、启动 official eval。V3 类 workflow 的结构化阶段多、单题时间长、远端 Docker 与模型代理都可能抖动，因此评测 Harness 的第一目标是把基础设施错误、空 patch、语义失败和官方评测结果拆开，避免一个可恢复的 transport failure 被写成模型能力失败。

当前 Harness 的基本原则是：模型生成阶段产生 patch；checkpoint 在语义状态之外保存容器工作区 diff；watchdog 根据状态决定继续生成、暂停等待、恢复现场或启动 official eval；official eval 完成后再把 resolved/unresolved 写回状态文件。任何真实 API key、临时代理 token 或上游 token 都不能进入日志、Markdown 或远端持久文件。

\section{整体结构}

\begin{center}
\begin{tabularx}{0.98\linewidth}{L{3.2cm}L{5.2cm}Y}
\toprule
层级 & 代表文件 & 职责 \\
\midrule
Workflow runtime & \code{opencollab/opencollab/application/workflow.py} & 管理结构化 agent、transport retry、failure counts、checkpoint 读写、worktree sidecar 捕获与恢复。 \\
Eval harness & \code{opencollab/opencollab/harness/evaluator.py} & 创建容器环境，注入 SWE-bench 测试，接入 checkpoint 恢复，捕获 budget、infra、timeout 和异常，返回可观测 EvalResult。 \\
Patch 生成器 & \code{swebench/gen_prediction_workflow.py} & 启动单题 workflow，抽取容器 diff，清理验证产物，空 patch 时从 checkpoint sidecar 抢救，写 predictions 与 metrics。 \\
并发启动器 & \code{scripts/swe_parallel_start.py} & 按任务拆分独立 run 目录和 tmux session，支持 \code{--resume}，控制远端代理模式、预算、超时和数据集。 \\
做题 watchdog & \code{scripts/swe_v3_wave_watchdog.py} & 汇总多个 run，分类任务状态，修复远端代理，触发 resume，发现可评测 patch 后启动 official eval。 \\
官方评测驱动 & \code{scripts/swe_auto_eval_driver.py} & 按 \code{record_id} 和 patch sha 创建 eval side-car，规范化 dataset，启动 official eval，发现 report，写 status。 \\
远端执行脚本 & \code{scripts/run_swe_v2_one_from_fifo.sh} & 在远端 tmux 内执行单题生成，接收受控 credential，设置代理环境，追加 generation log。 \\
\bottomrule
\end{tabularx}
\end{center}

\section{状态机}

watchdog 的状态分类是做题框架的核心。一个任务先从 \code{no_output} 或 \code{partial_output} 进入生成；生成中为 \code{generating}；若 patch 非空且 workflow done，则进入 \code{patch_done} 并等待 official eval；若 official eval 完成，则进入 \code{official_resolved} 或 \code{official_unresolved}。可恢复错误进入 \code{infra_paused}；已有 checkpoint sidecar 的中断任务进入 \code{recoverable_checkpoint}；空 patch 只在没有可恢复 sidecar 时进入 \code{empty_patch_invalid}。

\begin{longtable}{L{4.2cm}L{4.8cm}L{6.1cm}}
\toprule
状态 & 触发条件 & 下一步 \\
\midrule
\code{generating} & tmux session 仍活跃 & 只读监控，不启动 official eval。\\
\code{patch_done} & patch 非空且 \code{workflow_status=done} & 交给 \code{swe_auto_eval_driver.py} 启动 official eval。\\
\code{patch_done_advisory_gap} & patch 非空且 \code{workflow_status=advisory_gap} & 默认作为 \code{diagnostic_only} 候选保留；只有显式打开 \code{--eval-advisory-gap} 时才送 official eval。\\
\code{infra_paused} & transport、429、远端代理、Docker、tmux、SSH 等可恢复错误 & 等待健康恢复和冷却时间，随后 \code{--resume}。\\
\code{recoverable_checkpoint} & 任务 inactive，patch 为空，但 checkpoint sidecar 存在 & 新容器恢复 \code{checkpoint.worktree.patch} 后继续 workflow。\\
\code{empty_patch_invalid} & 已写 prediction，patch 为空，且没有可用 sidecar & 保留日志，进入人工诊断。\\
\code{official_unresolved} & official eval 完成且 patch 未解决 & 进入语义分析和下一轮算法修正。\\
\code{official_resolved} & official eval 完成且 patch 解决 & 计为通过样本。\\
\bottomrule
\end{longtable}

当前 summary 结束条件也按这个状态机调整。\code{infra_paused} 和 \code{recoverable_checkpoint} 不计入完成数；official eval 完成、终结性 infra invalid、诊断候选和没有可恢复 sidecar 的 empty patch invalid 进入最终 summary。generation technical failure 保留为待修复状态，修好脚本、镜像、dataset 或远端环境后继续推进。

\section{结构化调用与 transport 错误}

V2 的主要污染来自结构化 stage 失败。旧逻辑中，Connection error 可能被转成 \code{None}，再被 workflow fallback 解释成 verifier fail。当前 runtime 将 transport failure 单列，并给 schema-bound stage 加同 label 重试。Connection error、timeout、HTTP 408/409/429、5xx、proxy、tunnel、remote protocol 等都归入 transport 类；429 单独标为 \code{rate_limit_error}。重试耗尽后抛出 \code{WorkflowInfraError}，由 evaluator 转成带 \code{retryable_infra} 的 workflow result。

\begin{center}
\begin{tabularx}{0.96\linewidth}{L{4.0cm}L{4.8cm}Y}
\toprule
错误类型 & 当前分类 & 影响 \\
\midrule
Connection error / timeout & \code{transport_error} & 原地重试，耗尽后进入 retryable infra。\\
HTTP 429 & \code{rate_limit_error} & 保留 pause reason，等待冷却后 resume。\\
schema capture 缺失 & \code{schema_failure_count} & 不混入 transport 统计。\\
语义 blocker & \code{semantic_blocker_count} & 只有模型成功返回 verdict 后才进入语义状态。\\
认证、权限、上下文长度 & non-retryable & 记录为配置或输入问题，避免盲目重试。\\
\bottomrule
\end{tabularx}
\end{center}

\section{Checkpoint 与 worktree sidecar}

checkpoint 现在保存两类信息。第一类是语义状态，位于 \code{checkpoint.json}，包括 workflow stage、结构化产物、token offset 和 runtime failure counts。第二类是容器工作区 diff，位于 \code{checkpoint.worktree.patch}，配套元数据位于 \code{checkpoint.worktree.json}。diff 捕获使用临时 git index 执行 \code{git add -A} 和 \code{git diff --cached --binary}，并排除 SWE-bench 注入测试文件，避免把官方测试 patch 作为模型修改保存。

恢复时，evaluator 在新容器创建后立即调用 \code{restore_checkpoint_worktree()}。若当前 worktree 为空，则应用 sidecar patch；若当前 worktree 已经有 diff，则跳过 apply，保护已有现场。若 sidecar patch 无法应用，metrics 写入 \code{checkpoint_sidecar_unusable}，任务保留为诊断对象，不能提交文本 fallback 或空 patch。

\begin{longtable}{L{4.8cm}L{10.2cm}}
\toprule
机制 & 作用 \\
\midrule
\code{checkpoint.worktree.patch} & 保存模型已经写入容器但尚未 official eval 的源码 diff。\\
\code{checkpoint.worktree.json} & 保存 patch sha256、字节数、stage、HEAD、status short、\code{preserved_previous_patch} 和 \code{submission_eligible}。\\
空 diff 与 capture failure 保留旧 sidecar & 新 checkpoint 捕获为空或捕获命令失败时，保留已有可校验 sidecar，避免空 patch 或临时基础设施错误覆盖非空 patch；这类 sidecar 的 \code{submission_eligible=false}，只用于恢复现场。\\
恢复前探测 current diff & 当前容器已有修改时跳过 sidecar apply，避免重复应用。\\
runtime 异常 checkpoint & budget、infra、timeout 和普通异常都会触发 \code{_runtime} checkpoint，尽量保留当时 diff。\\
\bottomrule
\end{longtable}

\section{这次审核发现与修复}

这次独立审核由两个只读 subagent 完成。Chandrasekhar 检查 V3.1 workflow 协议，指出 ETV 状态枚举、final-gate fallback、coverage gate 和重新仲裁分支的风险。Descartes 检查评测 Harness，指出 checkpoint sidecar、ready eval 候选身份、technical eval failure 回读和 report 技术失败分类的风险。两路审查共同指向同一个结论：workflow checkpoint 与 sidecar 的方向正确，official eval 的 hash side-car 与 tmux/flock 设计也可用；需要继续补厚的是 watchdog 和评测 driver 对半写、错配和可恢复异常的处理。

第一个问题是 \code{checkpoint_summary()} 过去依赖 \code{checkpoint.json} 能被成功解析。一旦这个 JSON 在 NFS 或进程中断下损坏，旁边仍然存在的 \code{checkpoint.worktree.patch} 会被忽略。随后 watchdog 可能把任务分到 fresh start，\code{swe_parallel_start.py} 会归档 trajectory 目录，sidecar 也随之被移走。修复后，watchdog 会独立检查 \code{checkpoint.worktree.patch} 和 \code{checkpoint.worktree.json}；即使 \code{checkpoint.json} 损坏，只要 sidecar patch 非空，就标为 \code{checkpoint_reusable}，下一次启动强制走 \code{--resume}。

第二个问题是 \code{predictions.jsonl} 与 \code{metrics.jsonl} 过去分别 append，进程若在两次写入之间退出，watchdog 可能把新 prediction 与旧 metrics 配成一对。修复后，生成器为每次产出写入同一个 \code{record_id} 和 \code{patch_sha256}；watchdog 与并发启动器只接受同 \code{record_id} 的成对记录。若发现 prediction 已写而对应 metrics 尚未落盘，任务进入 \code{partial_output} 并等待下一轮或 resume。watchdog 启动 \code{swe_auto_eval_driver.py} 时传入当前 \code{record_id} 与 \code{patch_sha}，official eval driver 只读取这个候选 patch，避免历史 patch 被误评。

第三个问题是一次新的空 patch 可能遮住旧的非空 patch。Pro Lite 十题里出现过“前面生成过 diff，续跑后新容器干净，最终写出空 patch”的症状。修复后，最新 prediction/metrics 记录控制当前任务状态；旧的非空且 workflow done patch 只写入 \code{last_good_patch_*} 字段，作为人工抢救和恢复线索，不能替代当前行触发 official eval。runtime 侧也同步加厚：空 diff 捕获时若已有非空 \code{checkpoint.worktree.patch}，metadata 写为 \code{empty_worktree_preserved_previous_patch}；capture failure 时若旧 sidecar 仍可校验，metadata 写为 \code{failed_preserved_previous_patch}。这两类保留 sidecar 的 \code{submission_eligible=false}，恢复时可用，最终提交时不会直接变成 prediction。

第四个问题是 summary 与 status JSON 过去直接 \code{write_text}。如果进程在写文件中途退出，下一轮 \code{load_json()} 会读到半截 JSON 并返回空状态。修复后，watchdog 的 status、final summary，official eval 的 state、summary、status、输入 prediction 和 dataset 都改成临时文件写入后 \code{os.replace} 原子替换。中断时保留旧状态，临时文件可以被后续清理。

第五个问题是结束条件过早。旧的 \code{_done_enough()} 把 \code{generation_technical_failed} 和 \code{technical_eval_failed} 算作完成，official eval watch 也把缺 dataset、缺依赖、缺镜像这类 blocked 状态算作 idle。修复后，可继续的技术状态不会触发最终 summary；watchdog 下一轮仍能修脚本、补镜像、恢复代理或重启 eval。

\section{自动 official eval}

official eval 是做题完成的判定器。生成 patch 后，watchdog 只有在 patch 非空、workflow done、生成 tmux inactive、同题 eval 未完成且未活跃时，才把任务列入 \code{ready_for_eval}。若结果带 \code{advisory_gap} 或 \code{done_with_advisory_gap}，默认写成 \code{diagnostic_only}；显式打开 \code{--eval-advisory-gap} 才按诊断模式送评。eval driver 根据 run 目录、instance id、\code{record_id} 和 patch sha 创建唯一 side-car，避免同一题历史 patch 与新 patch 混用。

\code{swe_auto_eval_driver.py} 做三件事：发现 candidate patch、启动官方评测、读取 report。它支持 SWE-bench Lite 和 Swe Batch Pro Lite 数据集发现；遇到单对象 JSON 会规范化成单元素 dataset；report 查找覆盖 side-car 内 report、官方 harness 嵌套 report 和 \code{/home/dongyh} 下的全局 report。driver 现在按 \code{record_id} 和完整 patch sha 过滤候选，并反查同一 \code{record_id} 且同一 patch hash 的 metrics；没有 hash 一致的 eligible metrics，candidate 会写成 \code{blocked_missing_metric} 或 \code{blocked_metric_not_eval_eligible}，不会启动 official eval。report 中的 patch apply 失败、error ids、empty patch ids、缺 report、缺镜像、缺 dataset、缺依赖都写成技术评测状态，并与 official unresolved 分开。

\section{远端代理与 credential}

模型调用支持三种 credential 模式。常用模式是 \code{proxy-token} 或 \code{proxy-open}，远端通过 SSH 反向隧道访问本机代理。短期离线需求可以使用 \code{encrypted-upstream-env}，远端脚本用临时密钥解密 env，启动本地代理后删除明文 env 文件。无论哪种模式，日志脱敏和状态文件都不能输出真实 token。

远端代理 health 只能证明端口活着，无法证明每次上游调用都稳定。因此 Harness 不能把 health OK 当成结构化 stage 一定成功的证据。当前设计把代理 health 用于是否修复隧道、是否解除 \code{infra_paused} 的条件之一；实际结构化调用仍由 runtime 记录 transport failure、rate limit failure 和 retry attempts。

\section{测试与自检}

当前已有针对性测试覆盖了几类重要边界。这次使用 \code{py_compile} 覆盖改动脚本，并在 \code{opencollab/.venv} 中执行 \code{pytest}，目标文件 83 个测试全部通过。新增测试覆盖了损坏 checkpoint JSON 仍识别 sidecar、\code{record_id}+patch hash 双重绑定、最新空 patch 不会被旧非空 patch 替代、\code{last_good_patch_*} 仍保留抢救线索、blocked eval 等待修复、advisory gap 默认诊断、preserved sidecar 禁止最终提交、空 diff 和 capture failure 显式保留上一份可恢复 patch。

\begin{center}
\begin{tabularx}{0.96\linewidth}{L{4.2cm}Y}
\toprule
测试文件 & 覆盖内容 \\
\midrule
\code{opencollab/tests/test_workflow_context.py} & checkpoint token offset、diff command 注入路径排除、rate limit 分类。\\
\code{opencollab/tests/test_gen_prediction_workflow.py} & sidecar 恢复、当前 worktree 非空跳过、sidecar apply 失败时禁止提交 fallback、preserved sidecar 只可恢复不可提交。\\
\code{opencollab/tests/test_swe_eval_harness_scripts.py} & recoverable checkpoint 进入 resume group，retryable infra 进入暂停，Docker daemon failure 暂停，auto-eval 需要 eligible metrics。\\
\code{opencollab/tests/test_evaluator_workflow.py} & WorkflowInfraError 被转换成 \code{infra_invalid} result，checkpoint 写入 runtime 状态。\\
\bottomrule
\end{tabularx}
\end{center}

\section{当前保障边界}

当前 Harness 可以保障 checkpoint 边界之后的恢复。也就是说，只要 workflow 已经成功写过 checkpoint，语义状态、token offset 和容器 diff 都有机会在新容器里恢复。它也可以保障大多数可恢复基础设施错误进入暂停态，等待代理、Docker、SSH 或 tmux 恢复后继续。

当前 Harness 仍然无法覆盖任意时刻的物理中断。若进程在一次 LLM 调用未返回时被杀死，或者容器在 checkpoint 写入之前消失，这段调用消耗和未落盘修改仍可能丢失。若 sidecar patch 与新容器 base 不兼容，恢复会进入 \code{checkpoint_sidecar_unusable}，需要人工诊断。若 official eval 依赖的镜像或 dataset 缺失，自动脚本会阻塞并保留证据，但不会把 patch 判成 resolved 或 unresolved。

\section{独立审核结论}

两路审核的共同结论可以概括为：V3 失败暴露出的核心风险已经从 workflow 语义转移到 harness 可靠性。经过这次补厚，sidecar patch 在 checkpoint JSON 损坏、空 diff 捕获、capture failure 时仍能被发现或显式保留；prediction 与 metrics 有稳定配对键；official eval 只评当前 \code{record_id}/patch sha；状态与 summary JSON 采用原子替换；可恢复技术状态停在等待修复的状态，而非提前进入最终 summary。

这意味着之后再遇到模型代理断连、远端 Docker 异常、tmux 中断、official eval 缺镜像或 dataset 这类问题，自动脚本会尽量暂停并保留证据。修复环境后，watchdog 可以根据 checkpoint 和 sidecar 继续推进。若没有 checkpoint 或 sidecar，脚本会把任务保留为诊断对象，并避免把这种技术中断解释成 V3 算法失败。

\section{评审迭代门槛}

这次经验说明，复杂 workflow 和评测 Harness 不能以一次自检作为终点。之后所有 V3.1、V3.2 或做题脚本的结论都采用分级表述：一次本地 smoke 通过只表示语法和样例路径可运行；一次独立审查通过只表示该轮没有新增 P0/P1；只有连续两轮独立审查都没有新增 P0/P1，并且真实恢复 smoke 能从中断点继续生成非空 patch、按当前 \code{record_id}/patch sha 启动 official eval，才允许把该版本标为可进入小样本实验。

评审也要覆盖不同角度。一个审查者看 workflow 协议，重点检查 advisory gap、coverage gate、ETV failure、final gate fallback 是否被正确分类；另一个审查者看 Harness，重点检查 checkpoint sidecar、prediction/metrics 配对、状态 JSON 原子写、official eval report 分类和 retry/pause 语义。任何审查者提出 P0/P1 后，都要先修代码和测试，再重新编译报告。没有新 P0/P1 时，也要保留剩余风险，例如 LLM 调用未返回前的物理中断、sidecar 与新容器 base 不兼容、远端镜像缺失这类边界。

\section{下一步建议}

下一步应做一个真实恢复 smoke。选择 \code{django__django-11019} 或 Pro Lite 中已有 sidecar patch 的样本，启动 V3 生成后在 checkpoint sidecar 出现时人为断开远端代理或杀掉 tmux，让 watchdog 进入 \code{infra_paused} 或 \code{recoverable_checkpoint}。恢复代理后运行同一 watchdog，确认启动命令带 \code{--resume}，确认新容器日志出现 \code{checkpoint worktree restored}，确认最终 patch 非空并自动进入 official eval。

这个 smoke 通过后，V3.1 还应减少硬结构化 stage 数量。candidate contracts、post-validation factory、judge triage 这类辅助 stage 应改成 advisory evidence producer，连接失败时复用缓存或跳过；最终法庭保留为少数硬裁决。这样可以降低 APIConnectionError 对整题完成率的乘法伤害。

\end{document}
